\chapter{Flow - Graph Matching}

\section{Hall's Marriage Theorem}
    Description: Hall's marriage theorem is a fundamental result in combinatorics that provides a necessary and sufficient condition for a perfect matching to exist in a bipartite graph, or equivalently, for a transversal to exist for a collection of sets, often explained through matchmaking scenarios (e.g., can every girl find a compatible boy?). The core condition (Hall's Condition) states that for any subset of vertices (say, "girls") in one partition, the size of their combined neighbors (the "boys" they know) must be at least as large as the size of the subset itself. 

    Theorem: A bipartite graph $G=(G \cup B,E)$ has a matching that covers $G$ if and only if Hall's Condition holds: for every subset $S \subseteq G$, $|N(S)| \geq |S|$

\section{Project Selection Problem}
    Problem: "You are given a set of $N$ projects and $M$ machines. The $i$-th machine costs $q_i$. The i-th project yields $p_i$ revenue. Each project requires a set of machines. If multiple projects require the same machine, they can share the same one. Choose a set of machines to buy and projects to complete such that the sum of the revenues minus the sum of the costs is maximized."

    Solution: One simple way to model the above problem as a minimum cut is to create a graph with $N + M + 2$ vertices. Each vertex represents the source, the sink, a project or a machine. We will note the source as $S$, the sink as $T$, the $i$-th project as $P_i$ and the $i$-th machine as $M_i$. Then, add edges of weight $p_i$ between the source and the $i$-th project, edges of weight $q_i$ between the $i$-th machine and the sink, and edges of weight $+\inf$ between each project and each of its required machines.

    Answer: (Sum of all $P_i$) - (answer when we call max flow on the above graph).



\kactlimport{DinicFlow.h}
\kactlimport{MaxFlowMinCost.h}
\kactlimport{GraphMatching.h}
\kactlimport{GraphFastMatching.h}
\kactlimport{Hungarian.h}